\section{Support Vector Machine}

We intend to back up the humanness of the chain outputted by the simulated annealing algorithm with this binary-output method.

The main point of this method is to transform our data space into some other one. In such a way that we can separate the data points with a simple plane, a linear classifier.

We will represent each point with a $p$-dimensional vector, where $p$ is the number of data points. The function that gives us those coordinates is called the kernel.

\subsection{The kernel}

Our kernel, $k:\mathcal C\times\mathcal C\to \mathbb R$,  will be the Gaussian radial basis function: $$k(\textbf x_i, \textbf x_j)=\exp \left( 
    -\gamma \left\| \textbf x_i -\textbf x_j \right\|^2 
\right),\ \gamma>0,$$ where $\{\textbf x_i\}$ are the data points.

After applying the kernel we have a matrix $M\in\mathbb R^{p\times p}$, the Gram matrix, with the coordinates of our data points in the new space.

\subsection{Decision function}

Using the lagrangian multipliers method for minimization, we get that the decision function, $f:\mathcal C\to\mathbb R$:
$$
f(\textbf x) = b + \sum_{i=1}^p\lambda_i s_i k(\textbf x_i,\textbf x),
$$
where $\lambda_i$ are the lagrange multipliers obtained through said method, $s_i$ is the \textit{tag} of the $i$-th data point (may be $+1$, when human, or $-1$, when murine) and $b$ is the bias, obtained as follows:
$$
b=s_p-\sum_{i=1}^p \lambda_i s_i k(\textbf x_i, \textbf x_p),
$$
where the $p$-th data point is a support one, that is, the closest to the plane (or any of them) and such that $\lambda_p>0$.